\documentclass[a4paper,slidestop,xcolor=pst,dvips,blue]{beamer}

\usepackage{beamerthemesplit}
\usepackage[utf8]{inputenc}
\usepackage[spanish]{babel}
\usepackage{graphicx}
\usepackage{pstricks} % PSTricks package
\usepackage{setspace}
\usepackage{multirow}
\usepackage{listings}
\usepackage{pgfpages}
\usepackage{hyperref}
\usepackage{etoolbox}
\usepackage{epstopdf}

\makeatletter
\patchcmd{\beamer@sectionintoc}{\vskip1.5em}{\vskip0.5em}{}{}
\makeatother

\setbeamercovered{dynamic}
\setcounter{tocdepth}{2}
\setbeamercolor{frametitle}{fg=black,bg=white}
\setbeamercolor{section in toc shaded}{fg=black}
\setbeamercolor{section in toc}{fg=red}
\setbeamercolor{subsection in toc shaded}{fg=black}
\setbeamercolor{subsection in toc}{fg=red}
\setbeamerfont{section in toc}{size=\small}
\setbeamerfont{subsection in toc}{size=\small}
\setbeamertemplate{section in toc shaded}[default][99]
\setbeamertemplate{subsection in toc shaded}[default][99]

\AtBeginSection[]
{\begin{frame}[c]
  \frametitle{Índice}
	\tableofcontents[currentsection,
        sectionstyle=show/shaded,
        subsectionstyle=hide]
\end{frame}}

\AtBeginSubsection[]
{\begin{frame}[c]
	\frametitle{Índice}
	\tableofcontents[
  		currentsection,
  		sectionstyle=shaded/shaded,
  		currentsubsection,
  		subsectionstyle=show/shaded/hide
		]
\end{frame}}

\setbeamercolor{frametitle}{fg=black,bg=white}

\setbeamertemplate{frametitle}{
	\begin{centering}
		\insertframetitle
		\par
	\end{centering}
}

\usetheme[secheader]{Boadilla} 

\title[Ing. Requisitos]{Ingeniería de Requisitos}

\author[Pablo Sánchez]{\alert{Pablo Sánchez}}

\institute[I2E]{
		   Dpto. Ingeniería Informática y Electrónica \\
		   Universidad de Cantabria \\
		   Santander (Cantabria, España) \\
		   p.sanchez@unican.es
}

\date{}

\begin{document}

\begin{frame}[c]
	\titlepage
	\begin{columns}
		\column{0.50\linewidth}
			\centering
			\includegraphics[width=.28\textwidth,keepaspectratio=true]{images/istr.eps}
		\column{0.50\linewidth}
			\centering
			\includegraphics[width=.25\textwidth,keepaspectratio=true]{images/uc.eps}
    \end{columns}
\end{frame}

\section{Datos Administrativos}

\subsection{Profesorado}

\begin{frame}[c]
	\frametitle{Profesorado}
	\begin{center}
		\alert{Pablo S\'{a}nchez Barreiro}  \\
		\begin{small}
		Despacho 1069 \\
		Departamento de Ingeniería Informática y Electrónica \\
		p.sanchez@unican.es \\
		\end{small}
	\end{center}
\end{frame}


\section{Objetivos y Temario}

\subsection{Objetivos}

\begin{frame}[c]
	\frametitle{Objetivos de la Asignatura}
    \begin{enumerate}[<+->]
			\item Comprender la importancia y papel de la \emph{Ingeniería de Requisitos}.
            \item Saber diseñar y ejecutar procesos de \emph{Ingeniería de Requisitos}.
			\item Saber aplicar técnicas para la \emph{Captura de Requisitos}.
			\item Saber aplicar técnicas para la especificación de \emph{Requisitos Funcionales}.
			\item Comprender la naturaleza de los \emph{Requisitos No Funcionales}.
			\item Saber especificar algunos \emph{Requisitos No Funcionales}.
			\item Saber especificar de manera básica \emph{Procesos de Negocio}.
	 \end{enumerate}
\end{frame}

\subsection{Resultados de Aprendizaje}

\begin{frame}[c]
	\frametitle{Resultados de Aprendizaje}
    	\begin{enumerate}[<+->]
			\item El alumno conocerá, comprenderá y sabrá utilizar la terminología propia de la \emph{Ingeniería de Requisitos}.
            \item El alumno sabrá diseñar y ejecutar procesos de \emph{Captura de Requisitos}, empleando diferentes técnicas para la identificación de requisitos, sabiendo escoger la más adecuada en función de cada situación.
            \item El alumno sabrá especificar requisitos funcionales utilizando para ello \emph{objetivos}, \emph{casos de uso} e \emph{historias de usuario}.
            \item El alumno sabrá especificar requisitos no funcionales, en especial los requisitos relacionados con la \emph{seguridad}.
            \item El alumno sabrá modelar procesos de negocio utilizando  \emph{BPMN 2.0}.
		\end{enumerate}
\end{frame}

\subsection{Temario}

\begin{frame}[c]
	\frametitle{Temario}
	\begin{enumerate}[<+->]
		\item Procesos de Ingeniería de Requisitos.
		\item Técnicas de Captura de Requisitos.
		\item Especificación y Modelado de Requisitos Funcionales.
		\item Especificación y Modelado de Requisitos No Funcionales.
		\item Especificación y Modelado de Procesos de Negocio.
	\end{enumerate}
\end{frame}

\section{Metodología}

\subsection{Plataforma}

\begin{frame}[c]
	\frametitle{Plataforma de Trabajo}
	\begin{itemize}[<+->]
		\item Las plataformas de trabajo para la asignatura serán \emph{Moodle}, \emph{YouTube} y \emph{Discord}.
		\item \emph{Moodle} se utilizará para publicar material docente y gestionar entregas.
		\item \emph{YouTube} se utilizará para:
			\begin{itemize}[<+->] 
				\item Retransmitir y guardar algunas clases de aula.
				\item Visualización de vídeos en horas de trabajo autónomo.
			 \end{itemize}
		\item \emph{Discord} se utilizará para gestionar trabajo en grupo y tutorías.
		\item Todas las notificaciones se harán a través de \emph{Moodle} y/o \emph{Discord}.
		\item Es obligación del alumno estar atento a las posibles notificaciones que se realicen a través de dichas plataformas.
	\end{itemize}
\end{frame}

\subsection{Actividades}

\begin{frame}
	\frametitle{Clases de Aula}
	\begin{block}{Objetivo}
        Entender los conocimientos teóricos que constituyen la base de las habilidades, destrezas y competencias a adquirir al final de la asignatura.
	\end{block}
	\begin{itemize}
        \item<2-> Requerirán la preparación de material fuera del aula (\emph{flipped learning}).
		\item<3-> Resolución autónoma, dirigida y supervisada de ejercicios.
		\item<4-> Clases magistrales puntuales usando pizarra y/o transparencias.
		\item<5-> Las \alert{transparencias} son material de apoyo, \alert{no apuntes}.
		\item<6-> \alert{Sin conocimiento teórico no puede haber habilidades prácticas}.
	\end{itemize}
\end{frame}

\begin{frame}[c]
	\frametitle{Clases en Laboratorio}
	\begin{block}{Objetivo}
		Aplicar los conceptos teóricos aprendidos en las clases de aula al desarrollo de un sistema software real de mediana escala, con el objetivo
        de desarrollar las competencias procedimentales y actitudinales deseadas.
	\end{block}
	\begin{itemize}
        \item<2-> Desarrollo de un proyecto de especificación de requisitos.
		\item<3-> Dos etapas diferenciadas: \emph{definición} y \emph{ejecución} del proceso de Ingeniería de Requisitos.
        \item<4-> Hay etapas y calificaciones tanto colectivas como individuales.
        \item<5-> Los resultados se utilizarán para la realización del proyecto integrado de cuarto.
	\end{itemize}
\end{frame}

\subsection{Tutorías}

\begin{frame}[c]
	\frametitle{Horario Tutorías}
    \begin{itemize}[<+->]
        \item Preferentemente online, a través de Discord.
        \item Disponibilidad Martes y Jueves por las mañanas, salvo excepciones.
        \item Si se quiere asegurar disponibilidad, avisar con antelación.
	\end{itemize}
\end{frame}

\section{Métodos de Evaluación y Calificación}

\subsection{Fórmula Calificación}

\begin{frame}[c]
	\frametitle{Cálculo de la Calificación Final}
	\begin{block}{Fórmula de Cálculo de la Calificación Final}
		\begin{scriptsize}
        \begin{tabular}{ll}
			Calificación Final = & SI Proyecto $\geq$ 5.00 \\
                                 & \ \ AND Participación ENTONCES \\
                                 & \ \ \ SI Prueba Final  $\geq$ Proyecto ENTONCES \\
                                 & \ \ \ \ \ \ Prueba Final + 1 + Participación \\
                                 & \ \ \ SINO \\
                                 & \ \ \ \ \ \ Proyecto + Participación \\
                                 & \ \ \ FINSI \\
                                 & SINO \\
                                 & \ \ \ avg(Prueba Final, Proyecto)  \\
                                 & FINSI \\
		\end{tabular}
        \end{scriptsize}
	\end{block}
	\begin{itemize}
		\item<2-> Es necesario presentarse a la prueba final para superar la asignatura.
        \item<3-> La participación es un extra que variará entre 0 y 0.5.
        \item<4-> En caso de no haber participación, la calificación mínima de la prueba final es de 4.
        \item<5-> Una calificación media de 4.99 es suspenso.
	\end{itemize}
\end{frame}

\subsection{Métodos y Criterios de Evaluación}

\begin{frame}
	\frametitle{Prueba Evaluable Escrita Final}
	\begin{itemize}[<+->]
		\item Prueba evaluable escrita en las fechas asignadas por el centro.
		\item La prueba contendrá ejercicios y razonamientos teóricos.
		\item Tendrá una duración de unas 4 horas, dividida en dos partes.
        \item Se podrá hacer uso todo tipo de material escrito.
        \item En ningún caso se dejará hacer uso de dispositivos electrónicos, especialmente, de aquellos con capacidades de comunicación inalámbrica.
		\item \alert{El material escrito debe servir para consultar cuestiones puntuales}.
		\item Aquellas respuestas que sean una simple copia de lo disponible en algún tipo de material se calificarán con 0 puntos.
        \item Hay modelos de pruebas escritas disponibles en moodle.
	\end{itemize}
\end{frame}

\begin{frame}[c]
	\frametitle{Peso de las Etapas del Proyecto}
    \begin{center}
	\begin{tabular}{||l|r|c||}
    \hline \hline
    Actas y Planes de Trabajo                   &  5\% & G \\ \hline
    Definición de la Visión del Proyecto        & 10\% & G \\ \hline
    Definición del Contexto                     & 15\% & G \\ \hline
    Identificación de las Fuentes de Requisitos & 10\% & G \\ \hline
    Diseño del Plan de Captura de Requisitos    & 10\% & G \\ \hline \hline
    Diseño de una Actividad de Captura de Requisitos & 10\% & G \\ \hline
    Especificación de un Modelo de Objetivos    & 15\% & G \\ \hline
    Especificación de Requisitos Funcionales    & 15\% & I \\ \hline
    Especificación de Requisitos No Funcionales & 10\% & I \\ \hline \hline
    Modelado de Procesos de Negocio             &  5\% & I \\ \hline \hline
	\end{tabular}
\end{center}
\end{frame}

\begin{frame}[c]
	\frametitle{Evaluación del Proyecto}
	\begin{itemize}[<+->]
        \item Hay calificaciones tanto conjuntas como individuales en las actividades conjuntas.
        \item Calificación final determinada por la media ponderada de las diversas entregas.
        \item Hay que sacar un mínimo en la parte individual para poder superar la asignatura. 
	\end{itemize}
\end{frame}

\section{Bibliografía}

\begin{frame}[c]
	\frametitle{Bibliografía}
\begin{thebibliography}{}

\bibitem[\protect\astroncite{Pohl}{2010}]{pohl:2010}
Pohl, K. (2010).
\newblock {\em {Requirements Engineering: Fundamentals, Principles and
  Techniques}}.
\newblock Springer.

\bibitem[\protect\astroncite{Gray et~al.}{2012}]{gray:2012}
Gray, D., Brown, S., and Macanufo, J. (2012).
\newblock {\em {Gamestorming: 83 juegos para innovadores, inconformistas y
  generadores del cambio}}.
\newblock Deusto.

\bibitem[\protect\astroncite{van Lamsweerde}{2009}]{lamsweerde:2009}
van Lamsweerde, A. (2009).
\newblock {\em {Requirements Engineering: From System Goals to UML Models to
  Software Specifications}}.
\newblock Wiley.

\end{thebibliography}
\end{frame}

\begin{frame}[c]
	\frametitle{Bibliografía}
\begin{thebibliography}{}

\bibitem[\protect\astroncite{Cockburn}{2000}]{cockburn:2000}
Cockburn, A. (2000).
\newblock {\em {Writing Effective Use Cases}}.
\newblock Addison-Wesley.

\bibitem[\protect\astroncite{Cohn}{2004}]{Cohn2004}
Cohn, M. (2004).
\newblock {\em {User Stories Applied: For Agile Software Development}},
\newblock Addison-Wesley Professional.

\bibitem[\protect\astroncite{Chung et~al.}{1999}]{chung:1999}
Chung, L., Nixon, B.~A., Yu, E., and Mylopoulos, J. (1999).
\newblock {\em {Non-Functional Requirements in Software Engineering}}.
\newblock Kluwer Academic Puvlishers.

\end{thebibliography}

\end{frame}

\end{document}
